% GENERATED FILE. Edit canonical lessons, then run npm run generate:books.
% canonical-source-hash: fnv1a64:edc4906089d527b8
% canonical-lessons: TE-C258-paddhati, TE-C258-niyamam, TE-C258-prabhavam, TE-C258-abhivrddhi, TE-C258-paristhiti

\chapter{Method, Rule, Effect, Development, Situation}
\label{ch:te-method-rule-effect-development-situation}
\hlchaptermodality{258}

\begin{chapteropening}
\textbf{By the end of this chapter:} \emph{I can say a method, a rule, an effect, development and a situation.}

Complete the last lesson of chapter 258: I can say a method, a rule, an effect, development and a situation.
\end{chapteropening}

\section[\texorpdfstring{\te{పద్ధతి} (paddhati)}{పద్ధతి (paddhati)}]{\te{పద్ధతి} (paddhati) — a method, a way of doing things}

\label{lesson:TE-C258-paddhati}



\begin{quote}
Before the new one: say the Telugu for a quarrel, then the Telugu for culture.
\end{quote}

\subsection*{You'll want to know: \te{పద్ధతి}}
\textbf{\te{పద్ధతి}} — \emph{paddhati} — \textquotedblleft{}a method, a way of doing things\textquotedblright{}.

\begin{grammarlens}[title={What you've built}]
One more word for this chapter.
\end{grammarlens}

\subsection*{Your turn}
\begin{itemize}
\raggedright
  \item \emph{Say it:} \emph{paddhati}
  \item \emph{Say it:} \emph{paddhati}, once more
  \item \emph{Say it:} say \emph{paddhati}
  \item \emph{Recall:} say the Telugu for a quarrel, then the Telugu for culture, then say \emph{paddhati} again
\end{itemize}

\begin{tcolorbox}[breakable,colback=teal!4,colframe=teal!35!black,title={Before you move on}]
What does \te{పద్ధతి} mean? (\textquotedblleft{}A method, a way of doing things\textquotedblright{}.) Say it once more.
\end{tcolorbox}
\section[\texorpdfstring{\te{నియమం} (niyamaṁ)}{నియమం (niyamaṁ)}]{\te{నియమం} (niyamaṁ) — a rule}

\label{lesson:TE-C258-niyamam}



\begin{quote}
Before the new one: say the Telugu for culture, then the Telugu for a method.
\end{quote}

\subsection*{You'll want to know: \te{నియమం}}
\textbf{\te{నియమం}} — \emph{niyamaṁ} — \textquotedblleft{}a rule\textquotedblright{}.

\begin{grammarlens}[title={What you've built}]
One more word for this chapter.
\end{grammarlens}

\subsection*{Your turn}
\begin{itemize}
\raggedright
  \item \emph{Say it:} \emph{niyamaṁ}
  \item \emph{Say it:} \emph{niyamaṁ}, once more
  \item \emph{Say it:} say \emph{niyamaṁ}
  \item \emph{Recall:} say the Telugu for culture, then the Telugu for a method, then say \emph{niyamaṁ} again
\end{itemize}

\begin{tcolorbox}[breakable,colback=teal!4,colframe=teal!35!black,title={Before you move on}]
What does \te{నియమం} mean? (\textquotedblleft{}A rule\textquotedblright{}.) Say it once more.
\end{tcolorbox}
\section[\texorpdfstring{\te{ప్రభావం} (prabhāvaṁ)}{ప్రభావం (prabhāvaṁ)}]{\te{ప్రభావం} (prabhāvaṁ) — an effect, an influence}

\label{lesson:TE-C258-prabhavam}



\begin{quote}
Before the new one: say the Telugu for a method, then the Telugu for a rule.
\end{quote}

\subsection*{You'll want to know: \te{ప్రభావం}}
\textbf{\te{ప్రభావం}} — \emph{prabhāvaṁ} — \textquotedblleft{}an effect, an influence\textquotedblright{}.

\begin{grammarlens}[title={What you've built}]
One more word for this chapter.
\end{grammarlens}

\subsection*{Your turn}
\begin{itemize}
\raggedright
  \item \emph{Say it:} \emph{prabhāvaṁ}
  \item \emph{Say it:} \emph{prabhāvaṁ}, once more
  \item \emph{Say it:} say \emph{prabhāvaṁ}
  \item \emph{Recall:} say the Telugu for a method, then the Telugu for a rule, then say \emph{prabhāvaṁ} again
\end{itemize}

\begin{tcolorbox}[breakable,colback=teal!4,colframe=teal!35!black,title={Before you move on}]
What does \te{ప్రభావం} mean? (\textquotedblleft{}An effect, an influence\textquotedblright{}.) Say it once more.
\end{tcolorbox}
\section[\texorpdfstring{\te{అభివృద్ధి} (abhivṛddhi)}{అభివృద్ధి (abhivṛddhi)}]{\te{అభివృద్ధి} (abhivṛddhi) — development, progress}

\label{lesson:TE-C258-abhivrddhi}



\begin{quote}
Before the new one: say the Telugu for a rule, then the Telugu for an effect.
\end{quote}

\subsection*{You'll want to know: \te{అభివృద్ధి}}
\textbf{\te{అభివృద్ధి}} — \emph{abhivṛddhi} — \textquotedblleft{}development, progress\textquotedblright{}.

\begin{grammarlens}[title={What you've built}]
One more word for this chapter.
\end{grammarlens}

\subsection*{Your turn}
\begin{itemize}
\raggedright
  \item \emph{Say it:} \emph{abhivṛddhi}
  \item \emph{Say it:} \emph{abhivṛddhi}, once more
  \item \emph{Say it:} say \emph{abhivṛddhi}
  \item \emph{Recall:} say the Telugu for a rule, then the Telugu for an effect, then say \emph{abhivṛddhi} again
\end{itemize}

\begin{tcolorbox}[breakable,colback=teal!4,colframe=teal!35!black,title={Before you move on}]
What does \te{అభివృద్ధి} mean? (\textquotedblleft{}Development, progress\textquotedblright{}.) Say it once more.
\end{tcolorbox}
\section[\texorpdfstring{\te{పరిస్థితి} (paristhiti)}{పరిస్థితి (paristhiti)}]{\te{పరిస్థితి} (paristhiti) — a situation}

\label{lesson:TE-C258-paristhiti}



\begin{quote}
Before the new one: say the Telugu for an effect, then the Telugu for development.
\end{quote}

\subsection*{You'll want to know: \te{పరిస్థితి}}
\textbf{\te{పరిస్థితి}} — \emph{paristhiti} — \textquotedblleft{}a situation\textquotedblright{}.

\begin{grammarlens}[title={What you've built}]
One more word for this chapter.
\end{grammarlens}

\subsection*{Your turn}
\begin{itemize}
\raggedright
  \item \emph{Say it:} \emph{paristhiti}
  \item \emph{Say it:} \emph{paristhiti}, once more
  \item \emph{Say it:} say \emph{paristhiti}
  \item \emph{Recall:} say the Telugu for an effect, then the Telugu for development, then say \emph{paristhiti} again
\end{itemize}

\begin{tcolorbox}[breakable,colback=teal!4,colframe=teal!35!black,title={Before you move on}]
What does \te{పరిస్థితి} mean? (\textquotedblleft{}A situation\textquotedblright{}.) Say it once more.
\end{tcolorbox}
